\scanpage{16}

un scalaire sur aucune fibre résiduelle.

$A$ est appelé l'anneau des \uline{homothéties} (\inserted{sous-entendu :} \uline{directes}) \deleted{\unsure{de|pour}} pour la structure de similitude envisagée $\Sigma$, on le note $\mathcal A_\Sigma$.

\uline{Lemme 2} Soit $\Sigma$ une structure de similitude sur $E$, et $\mathcal A_\Sigma$ \inserted{$\subset\underline{\mathrm{End}}(E)$} l'Anneau des homothéties. Conditions équivalentes

a) $\deleted{\unsure{\Sigma|\Sigma_u}}\Sigma$ est strictement fidèle.

b) \inserted{$\mathcal A(s)$ est de rang 2 et} $\mathcal A(\deleted{\unsure{x}}s)-\underline{\mathrm{End}}(E(s))$ injectif pour tout $s\in X$, i.e. $\mathcal A$ sous-fibré vectoriel (de rang 2) de $\underline{\mathrm{End}}(E)$.

b$'$) (Si $\mathcal A=\mathcal A_u\overset{\mathrm{déf}}{=}\underline{\mathcal O}\cdot1_E+\underline{\mathcal O}\cdot u$) $\forall s\in X$, $u(s)\in\underline{\mathrm{End}}(E(s))$ n'est pas un scalaire

c) $E$ est un $\mathcal A$-Module \deleted{loc. libre de \unsure{rg|rang}} \inserted{inversible}, i.e loc. libre de rang 1.

c$'$) $E$ est un $\mathcal A$-module \inserted{loc. \deleted{\unsure{mon}}} monogène.

d) $\mathcal E=\underline{\mathrm{End}}(E)$ est un $\mathcal A$-Module \unsure{à|;} g. (pour la structure régulière gauche) loc libre de rang 2.

d$'$) $\mathcal E$ est un $\mathcal A$-module \unsure{loc\textsuperscript{t}|loc ?|loc.} engendré par 2 élém\textsuperscript{ts}.

\uline{Cor. 1}, Sous ces conditions, soit \deleted{$x\in\Gamma E$, alors $x$ est} \inserted{$\Delta\subset\underline{\mathrm{Quad}}(E)$} le sous-Module \deleted{loc. lib.} inv. de $\underline{\mathrm{Quad}}(E)$ tel que $\Sigma=\Sigma_\Delta$. On aura donc un iso. can.

\[
\begin{aligned}
&\Delta\simeq(\mathcal A/\mathcal O)\otimes L^\vee &&(\text{où } L=\Lambda^{\unsure{2|7}}E=\det E)\\
\text{d'où}\qquad&\Delta^\vee\simeq(\mathcal A/\mathcal O)^\vee\otimes L&&,
\end{aligned}
\]
% archive sheet number: 15
